\subsubsection{Kết quả chọn bộ theo dõi}

Bảng~\ref{tab:tracking-comparison} trình bày kết quả đánh giá hai thuật toán
theo dõi trên tập xác thực. BYTE đạt kết quả cao hơn SORT trên tất cả các
độ đo chính, đặc biệt ở AssA và IDF1, cho thấy khả năng duy trì danh tính đối
tượng tốt hơn trong quá trình theo dõi. Do đó, BYTE được lựa chọn cho các
thí nghiệm đánh giá trên tập kiểm tra ở các phần tiếp theo.

\begin{table}[htbp]
  \centering
  \caption{So sánh kết quả theo dõi của SORT và BYTE trên tập xác thực}
  \label{tab:tracking-comparison}
  \begin{tblr}
  {
    colspec = {lcc},
    vline{2,3}, hline{2},
  }
  Độ đo & YOLO + SORT & YOLO + BYTE \\
  \HOTA   & 68.455 & 69.633 \\
  \DetA   & 66.955 & 67.588 \\
  \AssA   & 70.343 & 72.072 \\
  LocA   & 87.653 & 87.669 \\
  \MOTA   & 73.978 & 75.607 \\
  \MOTP   & 86.294 & 86.311 \\
  \IDFone & 79.419 & 81.483 \\
  \end{tblr}
\end{table}

\subsubsection{Kết quả theo dõi trên tập kiểm tra}

Sau khi lựa chọn BYTE trên tập xác thực, cấu hình này được sử dụng để đánh giá
trên tập kiểm tra. Kết quả trong \tablename~\ref{tab:track-test-result} cho
thấy hệ thống đạt \HOTA\ bằng $60.716$ và \IDFone\ bằng $70.345$. Giá trị
\AssA\ đạt $62.800$ cho thấy khả năng duy trì danh tính đối tượng tương đối ổn
định, trong khi \DetA\ đạt $59.184$ phản ánh chất lượng phát hiện của mô hình
YOLO vẫn là yếu tố ảnh hưởng đáng kể đến hiệu năng theo dõi tổng thể. Ngoài ra,
các thông số về độ chuẩn xác định vị \MOTP\ ($84.339$) và LocA ($85.953$) cho
thấy các hộp bao được gán khá chuẩn xác với phương tiện.

\begin{table}[htbp]
  \centering
  \caption{Kết quả đánh giá quá trình theo dõi trên tập kiểm tra}
  \label{tab:track-test-result}
  \import{\tablespath}{track-test-result}
\end{table}

\figurename~\ref{fig:hota} cho thấy các độ đo \HOTA, \DetA\ và \AssA\ có xu
hướng giảm khi giá trị ngưỡng $\alpha$ tăng. Trong khoảng từ $\alpha=0.0$ đến
$\alpha=0.7$, mức suy giảm của các độ đo tương đối nhỏ, cho thấy hệ thống vẫn
duy trì được hiệu năng phát hiện và liên kết đối tượng ổn định ở các ngưỡng
đánh giá trung bình. Tuy nhiên, từ $\alpha \approx 0.7$ trở đi, tốc độ suy giảm
trở nên rõ rệt hơn, đặc biệt đối với \DetA.

\begin{figure}[htbp]
  \centering
  \includegraphics[width=0.7\textwidth]{hota}
  \caption{Kết quả đánh giá theo dõi bằng độ đo \HOTA}
  \label{fig:hota}
\end{figure}

So với \AssA, \DetA\ giảm nhanh hơn khi $\alpha$ tăng, cho thấy hiệu năng của
hệ thống bị ảnh hưởng nhiều hơn bởi chất lượng phát hiện đối tượng so với khả
năng duy trì danh tính. Nói cách khác, khi tiêu chuẩn đánh giá ngày càng nghiêm
ngặt, các sai lệch về vị trí hoặc kích thước hộp giới hạn trở thành nguyên nhân
chính làm giảm hiệu năng tổng thể.

Trong khi đó, LocA tăng dần theo $\alpha$ và đạt giá trị cao nhất ở vùng
$\alpha > 0.8$. Điều này cho thấy các đối tượng vẫn được xem là khớp ở các
ngưỡng cao thường có độ chính xác định vị tốt hơn. Do \HOTA\ được xác định từ
cả \DetA\ và \AssA, xu hướng giảm của \HOTA\ chủ yếu bị chi phối bởi sự suy
giảm của \DetA\ ở các giá trị $\alpha$ lớn.
